(a) Every nonempty open subset of an irreducible space is irreducible and therefore connected. A locally constant function on it is constant. Thus all restrictions between nonempty opens are the identity on the constant group, and restrictions to the empty set are also surjective.

(b) By \exref{II.1.8}, only surjectivity on sections remains. Fix $s\in\mathscr F''(U)$ and order pairs $(V,t)$, where $V\subset U$ is open and $t\in\mathscr F(V)$ lifts $s|_V$, by extension. A chain has an upper bound by gluing, so Zorn's lemma gives a maximal pair. If $V\ne U$, choose a local lift $t'\in\mathscr F(W)$ near a point outside $V$. On $V\cap W$, the difference $t'-t$ is a section of $\mathscr F'$. Flasqueness extends it to $W$. Subtracting this extension from $t'$ makes it agree with $t$, so they glue over $V\cup W$, contradicting maximality. Hence $V=U$.

(c) For $V\subset U$ and $s\in\mathscr F''(V)$, part (b) lifts $s$ to $\mathscr F(V)$. Flasqueness of $\mathscr F$ extends this lift to $U$, whose image extends $s$. Therefore $\mathscr F''$ is flasque.

(d) The restriction $(f_*\mathscr F)(U)\to(f_*\mathscr F)(V)$ is the restriction $\mathscr F(f^{-1}U)\to\mathscr F(f^{-1}V)$, hence is surjective.

(e) The groups $\mathscr G(U)=\prod_{P\in U}\mathscr F_P$ form a sheaf, since arbitrary functions glue uniquely. A section on $V\subset U$ extends by zero outside $V$, so $\mathscr G$ is flasque. The map $\mathscr F(U)\to\mathscr G(U)$, $s\mapsto(s_P)_{P\in U}$, respects restrictions and is injective, since a section whose germs all vanish is zero.
